% #############################################################################
% This is Chapter 4
% !TEX root = ../main.tex
% #############################################################################
\fancychapter{Methodology}
\label{chap:methodology}

This chapter describes the experimental methodology. \Cref{sec:method-models}
introduces the model family and the fine-tuning regimes.
\Cref{sec:method-task-formulations} explains how rewriting and identification
instances are represented in a shared sequence-to-sequence format.
\Cref{sec:method-training-stages} presents the three-stage fine-tuning
pipeline. \Cref{sec:method-ablation-design} defines the supervised protocol-selection
study preceding the 4B-parameter experiments.

\section{Models and Fine-Tuning Setup}
\label{sec:method-models}

The reported experiments use encoder-decoder models from the T5Gemma 2 family
\citep{zhang2025t5gemma2seeingreading} with 270M and 4B parameters. The 270M-parameter experiments start from \texttt{google/t5gemma-2-270m-270m},
and the 4B-parameter experiments start from \texttt{google/t5gemma-2-4b-4b}.

The 270M-parameter model is fully fine-tuned for the supervised
protocol-selection study described in \Cref{sec:method-ablation-design}. These
lower-cost experiments compare choices involving batch composition,
control-string treatment, and classification loss before the experiments with
the 4B-parameter model are conducted. The study is not a controlled
comparison between full fine-tuning and LoRA, since model scale and adaptation
method change at the same time.

The 4B-parameter model is adapted with LoRA using rank \(r=48\), scaling factor
\(\alpha=96\), and dropout \(=0.05\). LoRA is applied to the attention and
feed-forward projection layers, while bias terms remain frozen. With \(r=48\),
this setup yields 187,772,928 trainable LoRA parameters. A separate LoRA-rank
diagnostic is reported in \Cref{chap:appendix-lora-rank}.

The reported supervised experiments use seed 42. During supervised training,
encoder inputs and decoder targets are truncated to a maximum of 512 tokens
during preprocessing.

\section{Task Formulations and Instance Construction}
\label{sec:method-task-formulations}

Two control formulations are used to construct rewriting and identification
instances: encoder-controlled and decoder-controlled. Both use \texttt{PT} and
\texttt{BR}, while \texttt{CLS} appears only in encoder-controlled
identification inputs. \texttt{PT} marks an EP source sentence, \texttt{BR}
marks a BP source sentence, and \texttt{PT} and \texttt{BR} determine rewriting
direction and act as identification targets.

An aligned EP/BP pair can provide supervision for both tasks. For rewriting,
each sentence in the aligned pair is used once as input and once as output. One aligned pair therefore
produces a BP-to-EP instance and an EP-to-BP instance. For identification, each
sentence is treated independently and paired with its variety label. A
complete non-equal pair can consequently provide one BP identification
instance and one EP identification instance. Equal pairs, that is, aligned
EP/BP pairs whose two sentences are identical, are retained for rewriting,
but they do not generate identification rows. Such rows would assign different
variety labels to identical sentence text.
The templates used to build
these examples are summarized in \Cref{tab:method-training-templates}.
These templates correspond to the task-specific row types used during training.
\Cref{sec:method-stage-a} and \Cref{sec:method-stage-b} quantify the rewriting
and classification rows in each stage mixture.

In both formulations, rewriting and identification share the same
encoder-decoder parameters, embeddings, and decoder vocabulary. The
formulations differ in where the control signal is introduced and how
classification supervision is attached to the target sequence.

\subsection{Encoder-Controlled Formulation}
\label{sec:method-encoder-control}

For BP-to-EP rewriting, the encoder receives the BP source sentence prefixed
with \texttt{BR}, and the decoder target is the aligned EP sentence. For
EP-to-BP rewriting, the encoder receives the EP source sentence prefixed with
\texttt{PT}, and the decoder target is the aligned BP sentence. In the
encoder-controlled formulation, identification uses the input \texttt{CLS} followed
by the sentence, and the decoder target is the corresponding variety
label. In this task, the complete label target is supervised.

A matched rewriting-only configuration is defined as an ablation of this
standard encoder-controlled formulation. It uses the same rewriting templates
but omits identification examples during training.

\subsection{Decoder-Controlled Formulation}
\label{sec:method-decoder-label}

The alternative formulation moves the control signal to the generated
sequence. In rewriting instances, the encoder receives the clean source
sentence with no task prefix. For BP-to-EP rewriting, the decoder target starts
with \texttt{BR} and then continues with the EP rewrite. For EP-to-BP
rewriting, the decoder target starts with \texttt{PT} and then continues with
the BP rewrite. Identification examples use the same decoder-side labels but
omit the rewritten sentence, leaving \texttt{PT} or \texttt{BR} as the complete
target.

The coupling is particularly direct in the decoder-controlled formulation.
Both task types supervise the same first decoding decision, while rewriting
additionally requires generation of the converted sentence.

The decoder-controlled setup also requires a distinction in the loss mask.
Identification rows compute loss only for the first token of the tokenized
decoder target. \texttt{PT} and \texttt{BR} are already present in the
tokenizer's base vocabulary and each is represented as a single token in the
tested configurations, so this position corresponds to the complete variety
label. Rewriting rows retain the complete target sequence in the loss.
This masking isolates classification from sequence termination
and avoids teaching the decoder to treat the control string alone as a complete
rewriting output.

\begin{table}[htbp]
  \centering
  \thesistablefont
  \setlength{\tabcolsep}{4pt}
  \renewcommand{\arraystretch}{1.14}
  \caption{Sequence-to-sequence templates used to construct rewriting and identification instances. PT and BR denote the input variety, while CLS marks identification inputs.}
  \label{tab:method-training-templates}
  \begin{tabularx}{\textwidth}{@{}>{\hsize=.7\hsize}Y>{\hsize=1.15\hsize}Y>{\hsize=1.15\hsize}Y@{}}
    \toprule
    Setup & Encoder input & Decoder target \\
    \midrule
    Encoder-controlled rewriting &
    BP-to-EP: \texttt{BR} + BP sentence \newline EP-to-BP: \texttt{PT} + EP sentence &
    BP-to-EP: EP sentence \newline EP-to-BP: BP sentence \\
    \addlinespace[0.35em]
    Encoder-controlled identification &
    \texttt{CLS} + sentence &
    \texttt{PT} or \texttt{BR} \\
    \addlinespace[0.35em]
    Decoder-controlled rewriting &
    BP-to-EP: BP sentence \newline EP-to-BP: EP sentence &
    BP-to-EP: \texttt{BR} + EP sentence \newline EP-to-BP: \texttt{PT} + BP sentence \\
    \addlinespace[0.35em]
    Decoder-controlled identification &
    sentence &
    \texttt{PT} or \texttt{BR} \\
    \bottomrule
  \end{tabularx}
\end{table}

\section{Training Stages}
\label{sec:method-training-stages}

The training procedure is organized into three stages. Stage A provides initial
supervised adaptation, Stage B continues supervised training with targeted
resources, and Stage C optionally applies reward-based fine-tuning to the Stage
B models. The supervised protocol used in Stages A and B was selected beforehand
with the 270M-parameter model in the study described in
\Cref{sec:method-ablation-design}. Alongside this staged procedure, matched
variants are used to isolate selected training choices and are introduced in the
stage where they first apply. Both the 270M- and 4B-parameter models use Stages A
and B, whereas the Stage C configurations use the 4B-parameter model. Separate
experiments with the 270M-parameter model, reported in
\Cref{chap:appendix-rl-ablation}, probe isolated reward and candidate-group
design choices. These experiments serve as diagnostics rather than a selection
procedure for the 4B-parameter Stage C configuration.

Across the supervised stages, optimization uses AdamW
\citep{Loshchilov2017DecoupledWD} with weight decay \(0.01\). The learning rate
follows a linear schedule with warm-up, rising to the stage-specific peak value
and then decreasing over the remaining training steps. The
subsections give learning rates, warm-up budgets, gradient
accumulation, training-step budgets, data composition, and validation
procedures.

\begin{table}[htbp]
  \centering
  \thesistablefont
  \setlength{\tabcolsep}{5pt}
  \renewcommand{\arraystretch}{1.12}
  \caption{Task-specific training examples by stage and data condition. Percentages in parentheses give task shares within each training set.}
  \label{tab:method-training-composition}
  \begin{tabularx}{\textwidth}{@{}cYcc@{}}
    \toprule
    Stage & Data & Rewriting examples & Identification examples \\
    \midrule
    A & OpenSubtitles & 20,695,766 (73.0\%) & 7,666,883 (27.0\%) \\
    \midrule
    B & GPT-Wikipedia & 9,038 (57.6\%) & 6,652 (42.4\%) \\
    B & GPT-Wikipedia + FRMT & 14,042 (54.7\%) & 11,616 (45.3\%) \\
    \midrule
    C & GPT-Wikipedia & 6,686 (100\%) & --- \\
    C & GPT-Wikipedia + FRMT & 6,828 (100\%) & --- \\
    \bottomrule
  \end{tabularx}
\end{table}

\Cref{tab:method-training-composition} summarizes the task-specific examples
available for training at each stage. Stages A and B combine rewriting and
identification supervision, whereas Stage C uses rewriting-only training data.

\subsection{Stage A: Broad Variety Adaptation}
\label{sec:method-stage-a}

Stage A provides broad supervised adaptation to both Portuguese varieties using
OpenSubtitles as its only data source. OpenSubtitles provides substantially
larger-scale EP/BP supervision than the resources used in Stage B. The resulting
checkpoint starts the more targeted Stage B adaptation.

After task construction, Stage A follows the natural OpenSubtitles task mixture
summarized in \Cref{tab:method-training-composition}, with approximately 73\%
rewriting and 27\% identification supervision. Under the supervised protocol
selected in \Cref{sec:method-ablation-design}, the two tasks are not explicitly
rebalanced.

The 4B-parameter model uses batch size 1, gradient accumulation 32, a peak
learning rate of \(10^{-5}\), 400 warm-up steps, and a maximum of 2500 training
steps. The protocol-selection experiments with the 270M-parameter model introduced in
\Cref{sec:method-models} use batch size 1, gradient accumulation 16, a peak
learning rate of \(2\times10^{-5}\), 200 warm-up steps, and a fixed budget of
2500 training steps. If the complete training budget is reached, these settings
correspond to at most 80,000 training-example presentations for the 4B-parameter
model and 40,000 for the 270M-parameter model. Under this protocol, training
examples are shuffled in a seed-controlled order. These budgets cover
only a small fraction of the Stage A training pool, so examples do not repeat
before the maximum step budget is reached.

OpenSubtitles has no resource-level validation split. Stage A therefore uses a
fixed 200-instance validation view sampled with seed 42 from the processed
training data after task construction. The sampled instances remain eligible for
training, so the validation view is not a disjoint held-out split. For the
4B-parameter model, validation and checkpointing are performed every
200 steps with patience 5. The protocol-selection experiments with the 270M-parameter model validate and
checkpoint every 150 steps with the same patience. Early stopping
is based on validation loss, and the checkpoint with the lowest validation loss
is restored for subsequent training rather than retaining the final training
step.

Matched no-Stage-A configurations start at Stage B rather than from a Stage A
checkpoint to assess the initial adaptation stage's contribution. Stage B
settings and data composition are matched.

\subsection{Stage B: Targeted Variety Adaptation}
\label{sec:method-stage-b}

Stage B shifts training from broad subtitle supervision toward more targeted
EP/BP resources. It uses the reviewed GPT-Wikipedia resource described in
\Cref{sec:data-gpt-wikipedia} as its primary data source. A second Stage B
condition incorporates the complete FRMT training split, allowing the effect
of data composition to be evaluated while keeping the same task templates and
training procedure. The task composition of both Stage B conditions is
summarized in \Cref{tab:method-training-composition}. Adding FRMT increases the
amount of both rewriting and identification supervision and shifts the mixture
slightly toward identification examples.

The GPT-Wikipedia + FRMT data mix concatenates the
corresponding GPT-Wikipedia and FRMT training splits and, separately, their
validation splits. The mix uses no weighting or resampling.

The 4B-parameter model uses batch size 1, gradient accumulation 32, a peak
learning rate of \(10^{-5}\), 150 warm-up steps, and a maximum of 1200 training
steps. The protocol-selection experiments with the 270M-parameter model use batch size
1, gradient accumulation 16, a peak learning rate of \(2\times10^{-5}\), 150
warm-up steps, and a fixed budget of 1200 training steps. If the complete
training budget is reached, these settings correspond to at most 38,400
training-example presentations for the 4B-parameter model and 19,200 for the
270M-parameter model. For the 4B-parameter model, both Stage B training sets are
smaller than the 38,400-presentation budget, so examples can repeat across
shuffled epochs. For the 270M-parameter model, the 19,200-presentation budget
exceeds the GPT-Wikipedia training set but not the GPT-Wikipedia + FRMT data
mix.

Each Stage B condition uses a held-out validation split constructed with the
same task templates as its training split. Matched encoder-controlled and
decoder-controlled configurations within the same data condition use the same
underlying training and validation examples. For the 4B-parameter model, validation and checkpointing are performed every
100 steps with patience 3. The protocol-selection experiments with the 270M-parameter model also validate
and checkpoint every 100 steps, using patience 5. Early stopping is based on validation loss, and the checkpoint with
the lowest validation loss is restored.

A secondary decoder-controlled ablation removes equal source-reference
rewriting examples from the supervised data while retaining identification
examples. The corresponding Stage C variant uses the same continuation data,
with aggregate results reported in Appendix E.

\subsection{Stage C: Reward-Based Fine-Tuning}
\label{sec:method-stage-c}

Stage C applies reward-based continuation to the corresponding Stage B
configurations using more conservative rewriting-only data. The GPT-Wikipedia
configuration starts from the model trained with GPT-Wikipedia data, while the
GPT-Wikipedia + FRMT configuration starts from the model trained with the
corresponding data mix. Stage C is not applied to the no-Stage-A
configurations.

To select conservative rewriting examples for Stage C, several measures are
computed over each aligned EP/BP pair to quantify the extent and type of the
differences between the two sentences. A case-insensitive sequence alignment
over word and punctuation tokens is used to identify changed regions. A changed
region corresponds to one contiguous non-matching alignment block. The edit
proportion measures the number of changed word tokens relative to the larger
sentence, while structural similarity is computed from the sequence-alignment
ratio over the complete token sequences. Changes are also compared with a
predefined inventory of EP/BP marker forms. Changed word tokens in the
non-matching alignment blocks are divided into marker and non-marker tokens
according to case-insensitive membership in this inventory. The non-marker
change proportion is the number of non-marker changed-token occurrences across
both sentences divided by the larger sentence word count.

For a non-equal pair to be included in the Stage C training set, it must contain at
most four changed regions, have an edit proportion of at most 0.30, and have a
structural similarity score of at least 0.72. The non-marker change proportion
must not exceed 0.18, and the pair may contain at most six non-marker
changed-token occurrences. When at least one marker token changes, the number of
non-marker changed tokens may exceed the number of marker changed tokens by at
most two. This final constraint is not applied when no marker token changes.
Unchanged pairs are retained separately as no-change examples, for which copying
the source is the correct output.

Both Stage C training sets are balanced across rewriting directions. The
combined set contains the same 6,686 GPT-Wikipedia rewriting examples used in
the GPT-Wikipedia set together with 142 FRMT examples derived from 71 aligned
pairs. No-change examples make up about 35\% of each Stage C set.

Stage C contains no identification examples. Under decoder control, rewriting
targets still begin with \texttt{PT} or \texttt{BR}, and correctness of this
first generated token contributes to the reward.

The Stage C configurations with the 4B-parameter model use batch size 1,
gradient accumulation 8, a learning rate of \(5\times10^{-6}\), 20 warm-up
steps, and a fixed optimization budget of 600 steps. AdamW is used with weight decay 0 and the same linear learning-rate
schedule defined for the preceding stages. Unlike the supervised stages,
Stage C does not use validation-based early stopping or best-validation-loss
checkpoint selection. The configurations use seed 42 and gradient clipping with
maximum norm 1.0. The final training state after 600 steps is used for
evaluation.

The optimization uses a group-relative update inspired by the reward
optimization methods described in \Cref{sec:rl-reward-ft}
\citep{Shao2024DeepSeekMathPT}. For each training
input, four candidate rewrites are sampled with temperature 0.8, top-\(p\)
0.95, a maximum completion length of 128 tokens, and no n-gram blocking. The
gold target is then included in the same candidate group, so each group usually
contains five candidate entries. The gold target participates in reward
normalization like the sampled candidates.

The reward compares each rewrite with the reference and unchanged
source. For the BLEU component, let \(B_m\) denote candidate sentence-level BLEU
and \(B_c\) the BLEU from copying the source:

\[
R_{\mathrm{BLEU}}
=
\frac{B_m}{B_m+B_c+\epsilon}.
\]

Similarly, let \(T_m\) be candidate TER and \(T_c\) the copy-source TER:

\[
R_{\mathrm{TER}}
=
\frac{T_c+\epsilon}{T_m+T_c+2\epsilon}.
\]

The constant \(\epsilon=10^{-8}\) prevents division by zero.

These forms reward a candidate relative to the unchanged-source alternative
rather than using the raw BLEU and TER values alone. Encoder-controlled models
weight the two components equally:

\[
R_{\mathrm{enc}}
=
0.5R_{\mathrm{BLEU}}
+
0.5R_{\mathrm{TER}},
\]

Decoder-controlled models additionally reward correctness of the first
generated input-variety label:

\[
R_{\mathrm{dec}}
=
0.4R_{\mathrm{BLEU}}
+
0.4R_{\mathrm{TER}}
+
0.2R_{\mathrm{label}}.
\]

For BP-to-EP examples, \(R_{\mathrm{label}}\) is 1 when the first generated
label is \texttt{BR}. For EP-to-BP examples, it is 1 when the first generated
label is \texttt{PT}. Otherwise, it is 0. After removal of any decoder control
label and whitespace normalization, \(0.10\) is subtracted from the combined
candidate reward when the generated rewrite exactly matches the source and the
reference requires a change.

Candidate rewards are converted to group-normalized advantages using the
procedure described in \Cref{sec:rl-reward-ft}. If the reward standard deviation
within a group is zero, all advantages for that group are set to zero. The
implemented Stage C loss is

\[
\mathcal{L}_{\mathrm{StageC}}(\theta)
=
-\frac{1}{|G(x)|}
\sum_{y_i \in G(x)}
A_i
\frac{\log p_\theta(y_i \mid x)}{L_{\max}},
\]

where \(A_i\) is the group-normalized advantage for candidate \(y_i\) and
\(L_{\max}=128\) is the maximum completion length used for normalization. The
implemented objective does not use Proximal Policy Optimization (PPO)-style
likelihood ratios or clipping.
With \(\beta=0\), no explicit KL penalty toward the starting policy is applied.

Stage C differs from Stage B in both the optimization procedure and the
training data. Thus, Stage B-to-Stage C comparisons characterize the
Stage C procedure rather than the reward objective alone.

\section{Supervised Protocol Selection}
\label{sec:method-ablation-design}

The protocol-selection study with the 270M-parameter model compares four
decoder-controlled supervised-training configurations at lower computational
cost before the experiments with the 4B-parameter model. All configurations use full fine-tuning and follow the same Stage A and Stage
B training path described in \Cref{sec:method-training-stages}, with
OpenSubtitles used in Stage A and GPT-Wikipedia used in Stage B. The
comparison keeps the same T5Gemma 2 checkpoint, decoder-controlled task
formulation, task construction, evaluation procedure described in
\Cref{sec:evaluation-methodology}, and comparable step-budget settings, while varying only the supervised protocol component
being tested. The results are reported in
\Cref{sec:smaller-model-protocol-study}.

The comparison comprises one base protocol and three single-factor variants:
task-balanced training, added-token treatment, and restricted classification
loss. In the base protocol, rewriting and identification rows are combined in
the same training mixture and shuffled together, so their contribution follows
their relative frequency in the data rather than an explicit task balance. The
control strings use the existing base-vocabulary representations of
\texttt{PT} and \texttt{BR} described in \Cref{sec:method-decoder-label}.
Classification examples are scored against the normal decoder
vocabulary, with the supervised target positions determined by the formulation
described in \Cref{sec:method-task-formulations}.

The task-balanced variant changes only the composition of the examples
contributing to each optimizer update, enforcing equal numbers of rewriting
and identification rows. The data, task templates, tokenizer treatment, and
identification loss remain unchanged from the base protocol.

The added-token variant keeps the same control strings and task templates but
explicitly registers them with the
tokenizer instead of using their existing base-vocabulary representations.
The data mixture, sampling procedure, and loss computation remain unchanged
from the base protocol.

The restricted-classification-loss variant changes only classification loss by
restricting scoring to \texttt{PT} and \texttt{BR}. Rewriting examples still
use the full decoder vocabulary.
